\paragraph{Weight of paths at the segment cutoffs.}
\label{par:plane-typical}
\begin{lemma}[Weight of the paths obeying the cutoffs]
\label{lem:good-plane-path-count}
For segment \(j\) of Definition~\ref{hyp:coarsened-plane-frames}, let \(\mathcal T_j\) be the paths satisfying the degree cutoffs, with a choice of standard terminal tableau for the complement weight \(\kappa\). With \(D_k(b)=\binom{b+k-1}{k-1}\), their number satisfies
\begin{equation}\label{eq:plane-good-path-count}
 |\mathcal T_j|\ \ge\ (1-a^{-N_s^2})
    \prod_{x\text{ between frames}}D_k(n_x)
    \prod_{x\text{ incoming}}D_k(W_{j,x}-m_0).
\end{equation}
\end{lemma}

\begin{proof}
We prove this at the simultaneous cutoffs above. Let \(m\le N_s\)
be the number of letters in the segment and number the first
\(m-r\) steps before the outgoing frame. Use the hypotheses
\(k,\mathcal L\ge1088N_s^2\log a\). From the degree-loss bound
\(E'_x+E_x\le5b_x/(kY)\) we have \(a'_x\ge9b_x/10\).
For either kind of pre-outgoing row,
\[
 \frac{4a}{5P_{j,x}}\le B_x\le\frac{2a}{P_{j,x}},\qquad B_x\ge k.
\]
In fact the floor and padding lose at most
\[
 4kN_s\le b_x/(10P_{j,x}),
\]

since \(b_x/P_{j,x}\ge8Yk^2N_s\). Thus degrees of successive
rows of the same parity grow by at least
\[
 \rho=p^2/5>13.
\]

For any letter \(x\), each preceding pre-outgoing degree is at
most \(4\beta_x/p\), where \(\beta_x=a/P_{j,x}\), by strict
increase of the top positions.

First count all paths with the outgoing head conditions,
weighted by the number of tableaux of the terminal complement.
Before the outgoing frame the signed Pieri rules decompose the
tensor product of the indicated symmetric powers of the standard
module and its dual. Each path to a weight \(\lambda\) has weight
\(\dim V_\lambda\). After \(t\) outgoing steps remove the first \(t\)
entries, all equal to \(S\), and call the remaining row list the
tail.

At the next step its old and new tails \(\eta,\eta'\) obey
\[
 \eta_i\ge\eta'_i\ge\eta_{i+1}.
\]
Conversely every such
interlacing tail reinstates a vector strip on inserting the
extra head equal to \(S\), because all entries were at most
\(S\). Thus at the outgoing steps the ordinary branching rule
from \(\mathrm{GL}_{k-t}\) to \(\mathrm{GL}_{k-t-1}\) preserves the
total weight. These multiplicity-one rules and the Weyl dimension
formula apply to rational weights by a determinant twist and
duality. \cite{fulton-harris}
The last tail is the complement, whose tableau count is its
Weyl dimension. The total count before imposing the good path
inequalities is exactly
\[
 R=\prod_{x\text{ pre-outgoing}}D_k(B_x).
\]


\proofstep{The constraints on a path and its two signs.}
For a given path and sign \(\epsilon\in\{+,-\}\), let
\(\alpha^\epsilon\) be the nonincreasing list of positive parts
or of absolute negative parts of its tail. Up to the outgoing
frame the tail is the full list of \(k\) rows. Let \(\mathcal B^\epsilon\)
be the list of previously used pre-outgoing degrees of the
corresponding sign, with vectors positive and covectors negative.
For every \(T\ge0\) put
\[
 \delta_\epsilon(T)=\sum_{B\in\mathcal B^\epsilon}\min(B,T)
                    -\sum_h\min(\alpha^\epsilon_h,T).
\]
Each of these deficits starts at zero and is nondecreasing along
the path.

At a pre-outgoing step of the corresponding sign let
\(b\) be the strip size and \(\alpha,\alpha'\) the old and
new lists for this sign. Monotonicity and interlacing give
\[
 0\le\sum_h\bigl(\min(\alpha'_h,T)-\min(\alpha_h,T)\bigr)
           \le \min(b,T).
\]
The \(b\) bound sums the coordinate increments. For the
\(T\) bound use \(\alpha'_{h+1}\le\alpha_h\) and telescope. At a step of the opposite sign the
second sum can only decrease. During the outgoing frame both
lists of parts decrease entry by entry by the tail interlacing,
and the reference lists stay fixed.

Write
\[
 \Xi_x=\mathcal L\beta_x/(pk).
\]
We bound the weight of
paths violating any of the following conditions at a late top
\(x\): a covector decreases positive mass by at most \(\Xi_x\),
a vector decreases negative mass by at most \(\Xi_x\), and the
last plane entry before either step is at least \(2B_u/3\).
Here \(u\) is the earliest of the preceding \(r\) vectors, which
is before the outgoing frame. These conditions imply the good
path inequalities. At a covector all plane entries are
nonnegative, and the complement's positive mass cannot
increase, so the plane loss is at most the positive mass loss.

For a complement weight \(\kappa\), denote the positive and negative masses by \(m^+=\sum_i(\kappa_i)_+\) and \(m^-=\sum_i(-\kappa_i)_+\). At a vector step the change in the sum of the entries is 
\[
 \Delta\!\sum_i\kappa_i=\Delta m^+-\Delta m^-,
\]
 the positive mass gain plus the negative mass loss. Interlacing bounds its
positive mass gain by the last plane entry before the step.
Reading that entry backwards through the last \(r-1\) vector
steps bounds it by the first entry just after \(u\). Geometric growth of the vector degrees gives
\[
 \text{first entry just after }u\ \le\!\sum_{x\le u,\ x\text{ a vector}} B_x\le 2B_u.
\]
Using the bound
\[
 2B_u\le4a/P_{j,u},
\]
the complement
gain is at most
\[
       \text{negative mass loss}+2B_u
                      \le \Xi_x+4a/P_{j,u}.
\]

\proofstep{The loss from the mass test.}
For a violation of the mass condition let \(U\) be the sum
of preceding degrees of the opposite sign. If there is any
mass of that sign, at least one such pre-outgoing degree has
appeared. Its largest degree is at least \(k\) and at most
\(4\beta_x/p\). The geometric sum gives
\[
 U\le5\beta_x/p.
\]

The total mass of the sign before this step is at most \(U\)
by the deficit inequality at \(T\ge S\). Hence for \(T\ge U\)
the deficit of that sign grows by more than \(\Xi_x\). This
includes negative mass losses at outgoing steps, since removing
an entry equal to \(S\) preserves the negative mass.

\proofstep{The loss from the last-entry test.}
For a violation of the last-plane-entry condition, put \(B=B_u\).
If \(t\) outgoing steps have passed at this cut, \(u\) is the
oldest of the
\[
 h=r-t
\]
most recent pre-outgoing vectors.
The sum of positive rows after the first \(h\) in the tail is
at most the sum of all vector degrees strictly before \(u\).
At each of those \(h\) vector additions, interlacing bounds the
positive sum below row \(i+1\) after the addition by the sum
below row \(i\) just before. Covectors and the later outgoing
trims can only decrease these sums. Thus the positive sum after
row \(h\) is at most \(B/(\rho-1)\). Row \(h\) has positive part
less than \(2B/3\).

For
\[
 B\le T\le13B,
\]
the \(h\) latest pre-outgoing
vectors contribute at least \((h-1)T+B\) to the reference sum.
In the clipped tail, the first \(h-1\) rows contribute at
most \((h-1)T\), row \(h\) at most \(2B/3\), and the rest at
most \(B/(\rho-1)\). Subtracting gives
\[
 \delta_+(T)\ge \big((h-1)T+B\big)
                 -\big((h-1)T+2B/3+B/(\rho-1)\big)\ge B/4
\]
by the stipulated growth of the vector degrees.

\proofstep{Converting exceptional events into lost weight.}
To convert these deficits into a weight loss, set
\(J_k(T)=\sum_{l=1}^{k-1}(T+l)^{-2}\). For \(b\ge0\),
\[
 \log D_k(b)=\int_0^\infty\min(b,T)J_k(T)\,dT,
\]
where the left side for real \(b\) denotes the logarithm of the
binomial product. The equality follows by integrating each
\((T+l)^{-2}\) separately.

At the terminal complement weight
\(\lambda\) in dimension \(K\), the positive and negative lists
bound its dimension by
\[
 \dim V_\lambda\le\prod_h D_k(\alpha_h^+)D_k(\alpha_h^-).
\]
Indeed each Weyl factor is bounded by
\[
 1+\frac{\lambda_i-\lambda_j}{j-i}
 \le\left(1+\frac{\max(\lambda_i,0)}{j-i}\right)
     \left(1+\frac{\max(-\lambda_j,0)}{j-i}\right).
\]
Multiplying gives the bound with \(D_K\), which is at most
\(D_k\) for each nonnegative argument. Therefore
\(\log(R/\dim V_\lambda)\) is at least the integral of
\((\delta_++\delta_-)J_k\) at the last cut. If the mass bound
fails, this integral is at least
\[
 \Xi_x\sum_{l=1}^{k-1}\frac1{U+l}
 \ge \frac{\mathcal L(k-1)}{9k}\ge\frac{\mathcal L}{16},
\]
using \(k\le4\beta_x/p\). If the plane cutoff fails, it is at
least
\[
 \frac B4\sum_{l=1}^{k-1}
     \left(\frac1{B+l}-\frac1{13B+l}\right)
 \ge\frac{11(k-1)}{104}\ge\frac{k}{16}.
\]
Here \(B\ge k\), and both estimates are independent of the
remaining path and its strip degrees.

At each cut the two lists of nonzero parts have lengths at
most \(m\), since each pre-outgoing strip introduces at most
one new part of its sign. The outgoing interlacing keeps that
bound after removing the heads. Each part is an integer in
\([0,S]\). The lists padded to length \(m\), and the time of
the cut, determine the whole shape. There are at most
\((1+S)^{2m^2}\) shape paths because the zero start is fixed.
Thus the exceptional weight is at most
\[
 R\exp\bigl(2m^2\log(1+2aN_s)-\min(k,\mathcal L)/16\bigr)
 \le R a^{-N_s^2}.
\]
The last inequality uses \(1+2aN_s\le a^2\), supplied by
\eqref{eq:plane-count-size}, and the bounds on \(k,\mathcal L\).
This proves \eqref{eq:plane-good-path-count} for each choice
of intervening digits with the same constants.
\end{proof}
